\chapter{Rule-Aware Region Ranking (Phase 1.1)}
\label{chap:p1-1}

\section{Goal}

The pilot's single highest-leverage defect, identified in Chapter~10 of the pilot report and
restated as item~3.1 of the scale-up plan, is that candidate explanation regions were ranked by
pixel \emph{area}. Florence-2 returns one box for the worker and one box for the queried safety
object (a hard hat, a harness, a guardrail, an excavator); the worker's whole-body box is almost
always larger than a hard hat or a harness worn on that body, so area ranking put the \emph{worker}
at rank~1 in 96\% of PPE (rule\_1) samples. Every downstream masking and overlap metric then
perturbed or measured the worker box instead of the safety object the rule is actually about.

The pilot report converted this observation into a falsifiable prediction, which is the target this
phase must hit. Splitting the 163 pilot samples by which entity was ranked top-1, masking the worker
flipped the answer in only 37.5\% of cases versus 49.4\% for the actual safety object --- an
\textbf{11.9 percentage-point gap} that the report attributed entirely to the ranking heuristic
promoting the wrong entity, not to any property of the model. The prediction: once ranking stops
systematically promoting the worker, that gap should ``narrow substantially, if not close
outright.''

\section{What Was Done}

\subsection{A ranking policy, selected by flag, defaulting to the frozen behavior}

\texttt{regions.standardize\_regions} gained a \texttt{region\_ranking} parameter taking one of three
values:

\begin{itemize}
  \item \texttt{"area"} --- the frozen pilot behavior: model boxes sorted by descending area. This
  remains the default so that every recorded pilot number reproduces byte-for-byte.
  \item \texttt{"rule\_aware"} --- regions whose label is the rule's queried safety object are
  promoted above all others; within each tier, boxes are still ordered by descending area. This is
  the scale-up policy.
  \item \texttt{"attention"} --- reserved for Phase~4, where cross-attention / Integrated-Gradients
  mass will supply a real per-region importance signal. Until that signal exists, the mode raises
  \texttt{NotImplementedError} rather than silently falling back to area, so that no run can quietly
  produce area-ranked numbers while believing they are attention-ranked.
\end{itemize}

The two-tier sort is implemented as two stable sorts: first by area (descending), then by a tier key
(0 if the region's label is the rule's object, 1 otherwise). Python's stable sort preserves the area
ordering \emph{inside} each tier, so the result is exactly ``objects first, largest object first;
then everything else, largest first.'' \texttt{rule\_aware} requires a non-empty set of object
labels and raises \texttt{ValueError} if none is supplied, again refusing to degrade silently.

\subsection{One source of truth for the object-class map}

The rule$\rightarrow$object-label map (\{rule\_1: hard hat, rule\_2: harness, rule\_3: guardrail,
rule\_4: excavator\}) had been copy-pasted inline into three separate metric scripts
(\texttt{04\_extract\_regions}, \texttt{05\_descriptive\_accuracy}, \texttt{07\_stability\_test}).
It now lives once in \texttt{prompts.RULE\_OBJECT\_LABEL}, with a helper
\texttt{rule\_object\_labels(rule\_id)} returning the object-label set that \texttt{rule\_aware}
ranking consumes. The three scripts import the shared map instead of redeclaring it --- eliminating a
triplicated literal that was itself a latent instance of the hardcoded-assumption class of bug
Phase~1.7 targets.

\subsection{Configuration and record protection}

\texttt{config.REGION\_RANKING} holds the default policy (\texttt{"area"}). Scripts 04, 05, and 07
each gained a \texttt{-{}-region-ranking} command-line flag defaulting to that config value. Critically,
any non-default run writes to a \emph{mode-suffixed} output file
(\texttt{descriptive\_accuracy\_rule\_aware.csv}, not \texttt{descriptive\_accuracy.csv}), so a
scale-up validation run can never overwrite the frozen pilot record it is being compared against.

\subsection{Test-first}

Seven unit tests were written for the new behavior and watched to fail before any implementation
existed: object-over-larger-worker promotion, area ordering within each tier, grid fallback under
rule\_aware, the empty-\texttt{object\_labels} guard, the unknown-mode guard, and the
\texttt{attention} not-yet-implemented guard, plus a test pinning the default to area ranking.

\section{Results}

\subsection{The ranking now selects the right entity}

Re-ranking the 163 pilot samples' already-extracted boxes (no model re-inference needed for this
part) shows the policy doing exactly what it should. Among the 159 samples that had a model box at
all (four fell back to the grid under both policies):

\begin{center}
\begin{tabular}{p{6.2cm}p{3.2cm}p{3.2cm}}
\toprule
\textbf{Top-1 region is\ldots} & \textbf{area (frozen)} & \textbf{rule\_aware} \\
\midrule
the worker body          & 50.3\% & 0.6\% \\
the rule's queried object (rule\_1/2/3) & --- & 100\% \\
the rule's queried object (rule\_4)     & --- & 96\% \\
\bottomrule
\end{tabular}
\captionof{table}{Which entity the ranking promotes to rank~1, across the 159 model-box pilot
samples. Under area ranking the worker body was top-1 in just over half of all samples; under
rule\_aware it is top-1 in one.}
\label{tab:p11-entity}
\end{center}

For PPE (rule\_1) specifically --- the case the pilot cares most about, and the case the old ranking
handled worst --- the top region changed in 92.1\% of samples. It changed far less for rule\_3
(6.5\%) and rule\_4 (0\%), which is correct behavior, not a weakness: guardrails and excavators are
often already the largest box in frame, so area ranking and rule-aware ranking agree on them. The
fix moves precisely the samples that needed moving.

\subsection{The 11.9-point gap closes}

Reproducing the pilot report's worker-vs-object split on the actual re-run descriptive-accuracy
results:

\begin{center}
\begin{tabular}{p{4.4cm}p{4.4cm}p{4.4cm}}
\toprule
\textbf{} & \textbf{area (frozen)} & \textbf{rule\_aware} \\
\midrule
top-1 = worker: flip rate & 30.0\% ($n{=}80$) & 0.0\% ($n{=}1$) \\
top-1 = object: flip rate & 43.0\% ($n{=}79$) & 39.2\% ($n{=}158$) \\
worker-vs-object gap & $+13.0$ pp & eliminated (no worker group left) \\
\bottomrule
\end{tabular}
\captionof{table}{The worker-vs-object descriptive-accuracy gap, before and after. Under area
ranking, 80 of 159 samples masked the worker; under rule\_aware, one does. The systematic
gap is gone because the confound that produced it --- the ranking's inconsistency about which entity
it promotes --- is gone.}
\label{tab:p11-gap}
\end{center}

Two honest notes on these numbers. First, the frozen gap reproduces here as 13.0~pp using the raw
top-1 answer-flip signal, versus the report's 11.9~pp, which was computed on the bounded-completeness
``supported'' verdict (a signal that also folds in the top-2 mask). It is the same phenomenon
measured one step upstream; the magnitude agrees to within a point. Second, the object-group flip
rate drops slightly, from 43.0\% to 39.2\%, once ranking is fixed. This is not a regression: the 79
samples that previously masked the worker (mostly PPE) now join the object group, and their genuine
object-mask flip rate is lower than the rate of the samples that already had a large object on top.
The fix does not inflate the metric --- it de-contaminates it, and de-contamination reveals that
masking a small PPE object flips the proxy less often than the old, worker-contaminated aggregate
suggested. PPE's own descriptive accuracy nonetheless rises from 27.0\% to 33.3\% now that its top-1
mask lands on the hard hat rather than the person.

\subsection{The frozen record is intact}

Running \texttt{04\_extract\_regions.py -{}-region-ranking area} regenerated
\texttt{region\_extraction.csv} byte-for-byte identical to the committed pilot file (empty
\texttt{git diff}), confirming the refactor did not disturb the default path. The full unit suite is
green at 68 tests (61 original plus the 7 added here).

\section{A Problem This Surfaced, and Its Resolution}

Masking metrics carry a known artifact from the pilot: masking the \emph{worker} box can make the
worker undetectable on the rerun, which reroutes \texttt{\_answer\_presence\_rule} into a degenerate
scene-level fallback and flips the answer for a reason unrelated to the safety concept. The pilot
measured this ``worker-loss'' signature as 15 samples where masking top-1-plus-top-2 \emph{un-flips}
an answer that masking top-1 alone had flipped.

After rule-aware ranking, that count rose to 22, not fell. This looks alarming until traced: it is
the same artifact \emph{relocating}, not worsening. Under area ranking the worker was top-1, so
worker-loss struck at the top-1 mask; under rule-aware ranking the worker is now top-2, so worker-loss
strikes at the top-2 mask instead. The number of contaminated samples is similar; the stage at which
the contamination appears has moved because the worker moved down the ranking.

The resolution is architectural, not a patch to this phase: Phase~1.1's job is to make the metric
test the \emph{right box}, and it does. De-contaminating the worker-loss fallback is a separate,
explicitly-scoped item (Phase~1.4, worker-loss instrumentation), which will record post-mask worker
boxes and report each metric twice --- raw and worker-loss-corrected. This phase's result should
therefore be read as ``the ranking is now correct,'' and the aggregate descriptive-accuracy figure
should not be treated as trustworthy until Phase~1.4 lands. Keeping the two fixes as distinct items,
rather than entangling them, is what makes each independently verifiable.

\section{Standard vs. Non-Standard Choices}

\textbf{Standard.} Ranking explanation regions by semantic relevance to the query, rather than by
raw size, is ordinary practice; area was only ever a stopgap the pilot adopted because Florence-2's
grounding path exposes no native importance score.

\textbf{Non-standard, and the actual discipline.} Three choices go beyond the obvious fix. The old
behavior is preserved exactly, behind a default-valued flag, so the pilot record remains a
reproducible baseline rather than being overwritten by the improvement. Non-default runs are
physically prevented from clobbering that record by output-file suffixing. And the not-yet-built
\texttt{attention} mode fails loudly instead of silently falling back to area --- encoding the
principle that a metric must never quietly report a different computation than the one its
configuration claims. These are the habits that let a scale-up change a pipeline's behavior without
losing the ability to prove what the old behavior was.

\section{Implications}

Rule-aware ranking unblocks every masking and overlap metric (descriptive accuracy, sparsity,
stability, completeness) to operate on the safety object rather than the worker body, and it does so
with no additional model calls --- purely a re-ordering of boxes Florence-2 has already returned. The
pilot report's falsifiable prediction is met: the 11.9-point worker-vs-object gap closes outright.
The immediate caveat carried forward is that the worker-loss fallback, now relocated to the top-2
stage, must be instrumented (Phase~1.4) before the corrected descriptive-accuracy, robustness, and
completeness aggregates can be quoted.
